% ============================================================
\chapter{Characteristic Functions}
\label{ch:characteristic}
% ============================================================

\begin{intuition}
The characteristic function is a powerful analytic tool that transforms
probabilistic problems into problems of complex analysis. It always exists
(unlike the MGF), uniquely determines the distribution, and provides an elegant
route to the central limit theorem.
\end{intuition}

% ------------------------------------------------------------
\section{Definition and Basic Properties}
% ------------------------------------------------------------

\begin{definition}[Characteristic function]
The \textbf{characteristic function} of a real-valued random variable $X$ is
\[
    \varphi_X(t) = \E[e^{itX}] = \E[\cos(tX)] + i\,\E[\sin(tX)], \quad t \in \R.
\]
\end{definition}

\begin{theorem}[Fundamental properties]
\label{thm:cf_properties}
Let $\varphi_X$ be the characteristic function of $X$. Then:
\begin{enumerate}[label=(\roman*)]
    \item $\varphi_X(0) = 1$.
    \item $\abs{\varphi_X(t)} \leq 1$ for all $t \in \R$.
    \item $\varphi_X$ is \textbf{uniformly continuous} on $\R$.
    \item $\varphi_X(-t) = \overline{\varphi_X(t)}$ (complex conjugate).
    \item If $Y = aX + b$, then $\varphi_Y(t) = e^{ibt}\,\varphi_X(at)$.
    \item If $X \indep Y$, then $\varphi_{X+Y}(t) = \varphi_X(t)\,\varphi_Y(t)$.
\end{enumerate}
\end{theorem}

\begin{proof}[Proof of (iii)]
\begin{align*}
    \abs{\varphi_X(t+h) - \varphi_X(t)}
    &= \abs{\E[e^{itX}(e^{ihX}-1)]}
    \leq \E[\abs{e^{ihX}-1}].
\end{align*}
Since $\abs{e^{ihX}-1} \leq 2$ and $e^{ihX}-1 \to 0$ as $h \to 0$, dominated
convergence gives $\E[\abs{e^{ihX}-1}] \to 0$ uniformly in $t$.
\end{proof}

% ------------------------------------------------------------
\section{Moments from the Characteristic Function}
% ------------------------------------------------------------

\begin{theorem}[Moments and derivatives]
If $\E[\abs{X}^n] < \infty$, then $\varphi_X$ is $n$ times differentiable and
\[
    \varphi_X^{(k)}(0) = i^k\,\E[X^k], \quad k = 0, 1, \ldots, n.
\]
In particular:
$\E[X] = \varphi_X'(0)/i$ and $\E[X^2] = -\varphi_X''(0)$, whence
$\Var(X) = -\varphi_X''(0) - [\varphi_X'(0)/i]^2$.
\end{theorem}

% ------------------------------------------------------------
\section{Characteristic Functions of Standard Distributions}
% ------------------------------------------------------------

\begin{keyformulas}[title=\textbf{Table of characteristic functions}]
\begin{center}
\renewcommand{\arraystretch}{1.5}
\begin{tabular}{lc}
    \toprule
    \textbf{Distribution of $X$} & $\varphi_X(t)$ \\
    \midrule
    $\mathrm{Bernoulli}(p)$ & $1-p+pe^{it}$ \\
    $\mathrm{Bin}(n,p)$ & $(1-p+pe^{it})^n$ \\
    $\mathrm{Poisson}(\lambda)$ & $\exp(\lambda(e^{it}-1))$ \\
    $\mathrm{Geom}(p)$ & $\dfrac{pe^{it}}{1-(1-p)e^{it}}$ \\[4pt]
    $\mathcal{U}([a,b])$ & $\dfrac{e^{itb}-e^{ita}}{it(b-a)}$ \\[4pt]
    $\mathrm{Exp}(\lambda)$ & $\dfrac{\lambda}{\lambda-it}$ \\[4pt]
    $\mathcal{N}(\mu,\sigma^2)$ & $\exp\!\left(i\mu t - \dfrac{\sigma^2 t^2}{2}\right)$ \\[4pt]
    $\Gamma(\alpha,\lambda)$ & $\left(\dfrac{\lambda}{\lambda-it}\right)^\alpha$ \\[4pt]
    $\mathrm{Cauchy}(0,1)$ & $e^{-\abs{t}}$ \\
    \bottomrule
\end{tabular}
\end{center}
\end{keyformulas}

\begin{example}[Computation for the normal distribution]
If $X \sim \mathcal{N}(0,1)$:
\begin{align*}
    \varphi_X(t) &= \int_{-\infty}^{+\infty} e^{itx}\,\frac{1}{\sqrt{2\pi}}\,e^{-x^2/2}\,dx
    = \frac{1}{\sqrt{2\pi}}\int_{-\infty}^{+\infty} e^{-(x^2-2itx)/2}\,dx \\
    &= \frac{1}{\sqrt{2\pi}}\int_{-\infty}^{+\infty}
    e^{-(x-it)^2/2}\,e^{-t^2/2}\,dx = e^{-t^2/2}.
\end{align*}
For $X \sim \mathcal{N}(\mu, \sigma^2)$, write $X = \sigma Z + \mu$, so
$\varphi_X(t) = e^{i\mu t}\,e^{-\sigma^2 t^2/2} = e^{i\mu t - \sigma^2 t^2/2}$.
\end{example}

% ------------------------------------------------------------
\section{Uniqueness and Inversion Theorems}
% ------------------------------------------------------------

\begin{theorem}[Uniqueness theorem]
\label{thm:uniqueness}
Two random variables $X$ and $Y$ have the same distribution if and only if
$\varphi_X(t) = \varphi_Y(t)$ for all $t \in \R$.
\end{theorem}

\begin{theorem}[Inversion formula]
\label{thm:inversion}
If $\varphi_X \in L^1(\R)$ (i.e.\ $\int\abs{\varphi_X(t)}\,dt < \infty$),
then $X$ has a continuous density given by
\[
    f_X(x) = \frac{1}{2\pi}\int_{-\infty}^{+\infty} e^{-itx}\,\varphi_X(t)\,dt.
\]
\end{theorem}

\begin{remark}
The inversion formula is the probabilistic analogue of the inverse Fourier
transform. The characteristic function is simply the Fourier transform of the
distribution of $X$.
\end{remark}

% ------------------------------------------------------------
\section{L\'evy's Continuity Theorem}
% ------------------------------------------------------------

\begin{theorem}[L\'evy's continuity theorem]
\label{thm:levy}
Let $(X_n)$ be a sequence of random variables with $\varphi_n = \varphi_{X_n}$.
\begin{enumerate}[label=(\roman*)]
    \item If $X_n \xrightarrow{\mathcal{L}} X$, then $\varphi_n(t) \to \varphi_X(t)$
    for all $t$.
    \item Conversely, if $\varphi_n(t) \to \psi(t)$ for all $t$ and $\psi$ is
    \textbf{continuous at $0$}, then $\psi$ is the characteristic function of some
    random variable $X$ and $X_n \xrightarrow{\mathcal{L}} X$.
\end{enumerate}
\end{theorem}

\begin{remark}
This theorem is used in the proof of the CLT (Chapter~10): one shows pointwise
convergence $\varphi_{Z_n}(t) \to e^{-t^2/2}$, and continuity at $0$ of the limit
is obvious.
\end{remark}

% ------------------------------------------------------------
\section{Application: Stability of the Normal Distribution}
% ------------------------------------------------------------

\begin{theorem}[Stability of the normal]
If $X_1 \sim \mathcal{N}(\mu_1, \sigma_1^2)$ and
$X_2 \sim \mathcal{N}(\mu_2, \sigma_2^2)$ are independent, then
\[
    X_1 + X_2 \sim \mathcal{N}(\mu_1+\mu_2,\; \sigma_1^2+\sigma_2^2).
\]
\end{theorem}

\begin{proof}
\begin{align*}
    \varphi_{X_1+X_2}(t) &= \varphi_{X_1}(t)\,\varphi_{X_2}(t) \\
    &= e^{i\mu_1 t - \sigma_1^2 t^2/2}\,e^{i\mu_2 t - \sigma_2^2 t^2/2} \\
    &= e^{i(\mu_1+\mu_2)t - (\sigma_1^2+\sigma_2^2)t^2/2},
\end{align*}
which is the CF of $\mathcal{N}(\mu_1+\mu_2, \sigma_1^2+\sigma_2^2)$. The result
follows by uniqueness.
\end{proof}

% ------------------------------------------------------------
\section{Application: Stability of the Poisson Distribution}
% ------------------------------------------------------------

\begin{proposition}
If $X_1 \sim \mathrm{Poisson}(\lambda_1)$ and
$X_2 \sim \mathrm{Poisson}(\lambda_2)$ are independent, then
$X_1 + X_2 \sim \mathrm{Poisson}(\lambda_1+\lambda_2)$.
\end{proposition}

\begin{proof}
$\varphi_{X_1+X_2}(t) = e^{\lambda_1(e^{it}-1)}\cdot e^{\lambda_2(e^{it}-1)}
= e^{(\lambda_1+\lambda_2)(e^{it}-1)}$, the CF of
$\mathrm{Poisson}(\lambda_1+\lambda_2)$.
\end{proof}

% ------------------------------------------------------------
\section{Characteristic Functions and Convergence in Distribution}
% ------------------------------------------------------------

\begin{example}[Poisson to normal convergence]
Let $X_n \sim \mathrm{Poisson}(n)$. Set $Z_n = (X_n-n)/\sqrt{n}$.
\begin{align*}
    \varphi_{Z_n}(t) &= e^{-it\sqrt{n}}\,\varphi_{X_n}(t/\sqrt{n})
    = e^{-it\sqrt{n}}\,e^{n(e^{it/\sqrt{n}}-1)} \\
    &= \exp\!\left(n\!\left(e^{it/\sqrt{n}} - 1 - \frac{it}{\sqrt{n}}\right)\right).
\end{align*}
Taylor expansion: $e^{it/\sqrt{n}} = 1 + it/\sqrt{n} - t^2/(2n) + o(1/n)$, so
\[
    \varphi_{Z_n}(t) = \exp\!\left(n\!\left(-\frac{t^2}{2n}+o(1/n)\right)\right)
    \to e^{-t^2/2}.
\]
By L\'evy's theorem, $Z_n \xrightarrow{\mathcal{L}} \mathcal{N}(0,1)$.
\end{example}

% ------------------------------------------------------------
\section{Exercises}
% ------------------------------------------------------------

\begin{exercise}
Compute the characteristic function of $\mathcal{U}(\{-1, 0, 1\})$.
\end{exercise}

\begin{exercise}
Let $X \sim \mathrm{Exp}(\lambda)$. Compute $\varphi_X(t)$ and deduce
$\E[X]$, $\E[X^2]$, $\Var(X)$.
\end{exercise}

\begin{exercise}
Show that $\varphi_X(t) = e^{-\abs{t}}$ is the characteristic function of the
standard Cauchy distribution. Deduce that if $X_1, X_2$ are i.i.d.\ standard
Cauchy, then $(X_1+X_2)/2$ is again standard Cauchy.
\end{exercise}

\begin{exercise}
Using the inversion formula, recover the density of $\mathcal{N}(0,1)$ from its
CF $\varphi(t) = e^{-t^2/2}$.
\end{exercise}

\begin{exercise}
Show that the CF of $\Gamma(\alpha, \lambda)$ is
$(\lambda/(\lambda-it))^\alpha$. Deduce that the sum of $n$ i.i.d.\
$\mathrm{Exp}(\lambda)$ variables follows a $\Gamma(n, \lambda)$ distribution.
\end{exercise}

\begin{exercise}
Let $X_n \sim \mathrm{Bin}(n, \lambda/n)$ with $\lambda > 0$ fixed. By computing
$\varphi_{X_n}(t)$, recover the Poisson approximation:
$X_n \xrightarrow{\mathcal{L}} \mathrm{Poisson}(\lambda)$.
\end{exercise}

\begin{exercise}
Show that if $\varphi_X$ is real-valued, then $X$ and $-X$ have the same
distribution (the distribution of $X$ is symmetric about $0$).
\end{exercise}

\begin{exercise}
Let $(X_n)$ be i.i.d.\ with $\varphi_{X_1}(t) = 1/(1+t^2)$ (Cauchy). Show that
$\bar{X}_n$ has the same distribution as $X_1$ and explain why the CLT does not
apply.
\end{exercise}
